\documentclass[11pt]{amsart}
\usepackage{amssymb}
\usepackage[margin=1in]{geometry}
\usepackage[hidelinks]{hyperref}
\hypersetup{
  pdftitle={Post-Lie nonexistence with a semisimple target},
  pdfauthor={Alec Kriebel},
  pdfsubject={Finite-dimensional complex post-Lie algebras}
}
\newtheorem{theorem}{Theorem}
\newtheorem{corollary}[theorem]{Corollary}
\newtheorem{lemma}[theorem]{Lemma}
\theoremstyle{remark}
\newtheorem{remark}[theorem]{Remark}
\newcommand{\g}{\mathfrak g}
\newcommand{\n}{\mathfrak n}
\newcommand{\s}{\mathfrak s}
\newcommand{\rad}{\mathfrak r}
\newcommand{\C}{\mathbb C}
\newcommand{\Z}{\mathbb Z}
\DeclareMathOperator{\Ad}{Ad}
\DeclareMathOperator{\ad}{ad}
\DeclareMathOperator{\tr}{tr}
\DeclareMathOperator{\Der}{Der}

\title{Post-Lie nonexistence with a semisimple target}
\author{Alec Kriebel}
\address{Independent researcher}
\urladdr{https://orcid.org/0009-0001-9320-500X}
\date{16 October 2026; version 1.0.0}
\keywords{Post-Lie algebra, unimodular Lie algebra, semisimple Lie algebra,
holomorphic action}

\begin{document}
\begin{abstract}
Let \(\g\) be a finite-dimensional complex unimodular Lie algebra and
\(\n\) a complex semisimple Lie algebra of the same dimension.
We prove that if two Lie algebra homomorphisms
\(j_1,j_2:\g\to\n\) have invertible difference, then \(\g\) is
semisimple.  The integrated left-right action has a holomorphic
infinitesimal determinant.  Unimodularity makes this determinant
invariant, and its constancy on an open orbit extends to the whole
target group.  Transitivity and simple connectedness then identify the
source and target manifolds; their compact-form homology rules out a
nonzero source radical.  Consequently a perfect nonsemisimple complex
source and a semisimple target admit no post-Lie algebra structure.
\end{abstract}
\maketitle

\section{Introduction}

All Lie algebras in this note are finite-dimensional over \(\C\).
For two brackets on a common vector space, a post-Lie product relates
the source Lie algebra \(\g\) to the target Lie algebra \(\n\).
Burde, reporting joint work with Dekimpe, formulated the nonexistence
conjecture for a perfect nonsemisimple source and an arbitrary semisimple
target in Oberwolfach Report 46/2023, printed p.~2689
\cite{OWR}.  Burde--Dekimpe--Monadjem proved the target
\(\mathfrak{sl}_3(\C)\) case and further special cases
\cite[Theorem~3.10 and Propositions~3.11--3.12]{BDM}.
An earlier result of Burde--Dekimpe excludes a \emph{solvable
unimodular} source with a nonzero semisimple target
\cite[Theorem~3.3]{BD}.

We prove the conjectured nonexistence statement as a consequence of a
stronger obstruction for unimodular sources.  The passage from post-Lie
products to transverse pairs of homomorphisms is established in the
literature \cite[Proposition~2.16]{BDV},
\cite[Proposition~2.6]{BDM}; we include the reduction with explicit
bracket conventions.

The determinant method has a classical affine-space predecessor.
Baues associates a determinant relative invariant to an
\'etale affine representation and proves that triviality of its
character forces global transitivity
\cite[\S5.4.1, equation~(5.3), and Theorem~5.19]{Baues};
see also his Corollary~5.20 for the unimodular, volume-preserving case.
We use that method on a complex group: the determinant is globally
holomorphic, and the identity theorem replaces polynomial continuation
on affine space.  The subsequent obstruction uses classical
integration, Levi decomposition, compact real forms and ordinary
singular homology.  In particular, we do not invoke an affine-space
completeness theorem for the group target.

The literature comparison underlying this note is bounded: no prior
full theorem under the matching hypotheses was located in the primary
passages inspected.  This is not a claim of firstness, exhaustive
coverage, or current openness.  In particular, the original body of
Onishchik's 1969 work on reductive-group decompositions was inaccessible
in that comparison; this is an access limitation, not an identified
competing theorem.

\section{A unimodular transverse-pair obstruction}

Recall that \(\g\) is \emph{unimodular} if
\(\tr(\ad_{\g}x)=0\) for every \(x\in\g\).

\begin{theorem}\label{thm:pair}
Let \(\g\) and \(\n\) be complex Lie algebras of the same finite
dimension, with \(\g\) unimodular and \(\n\) semisimple.
Suppose \(j_1,j_2:\g\to\n\) are Lie algebra homomorphisms and
\(j_1-j_2\) is a vector-space isomorphism.  Then \(\g\) is semisimple.
\end{theorem}

We first recall the solvable-group fact needed in the proof
\cite[Theorem~49.1]{Etingof}, including its short induction.

\begin{lemma}\label{lem:solvable}
The connected simply connected complex Lie group with a solvable
Lie algebra \(\rad\) is diffeomorphic to
\(\C^{\dim_{\C}\rad}\), and is therefore contractible.
\end{lemma}

\begin{proof}
Induct on \(\dim_{\C}\rad\).  For \(\rad\ne0\), solvability gives a
nonzero functional vanishing on \([\rad,\rad]\).  Its kernel
\(\rad'\) is a codimension-one ideal.  A complementary line \(\C x\)
is a subalgebra, so \(\rad=\rad'\rtimes\C x\).
Let \(R'\) be the simply connected integration of \(\rad'\).
The automorphisms \(\exp(t\ad x)|_{\rad'}\), for \(t\in\C\),
integrate uniquely to automorphisms of \(R'\).  They give a
holomorphic action of the additive group \(\C\).
Thus \(R'\rtimes\C\) is simply connected and integrates \(\rad\).
Uniqueness of simply connected integration identifies it with the
required group, whose manifold is \(R'\times\C\).
Induction proves the assertion.  We use integration and uniqueness in
their complex, finite-dimensional form
\cite[Theorems~9.12--9.13 and Corollary~9.14]{Etingof}.
\end{proof}

\begin{proof}[Proof of Theorem~\ref{thm:pair}]
Put \(d=\dim_{\C}\g=\dim_{\C}\n\).
If \(d=0\), both algebras are zero and the conclusion holds.
Assume \(d>0\).
Choose connected simply connected complex Lie groups \(G,N\)
integrating \(\g,\n\).  The maps \(j_i\) integrate to holomorphic
homomorphisms \(J_i:G\to N\)
\cite[Theorems~9.12--9.13 and Corollary~9.14]{Etingof}.
They define a global holomorphic action
\begin{equation}\label{eq:action}
 A_h(p)=J_1(h)\,p\,J_2(h)^{-1}\qquad(h\in G,\ p\in N).
\end{equation}
No closedness or injectivity of the images \(J_i(G)\) is required.

Under left trivialization of \(TN\), the differential at the identity
of \(h\mapsto A_h(p)\) is
\begin{equation}\label{eq:B}
 B_p:\g\longrightarrow\n,\qquad
 B_px=\Ad_N(p^{-1})j_1x-j_2x .
\end{equation}
Fix bases of \(\g,\n\) and define the globally holomorphic function
\(f(p)=\det B_p\).  Since \(B_e=j_1-j_2\) is invertible, the
holomorphic inverse function theorem makes the orbit of \(e\) open.

For \(q=A_h(p)\), equivariance of the integrated homomorphisms gives
\begin{equation}\label{eq:covariance}
 B_q=\Ad_N(J_2(h))\,B_p\,\Ad_G(h^{-1}).
\end{equation}
Indeed, the first term on the right is
\[
 \Ad_N(J_2(h))\Ad_N(p^{-1})\Ad_N(J_1(h)^{-1})j_1,
\]
and the second is
\(-\Ad_N(J_2(h))j_2\Ad_G(h^{-1})=-j_2\).
These are the two terms of \(B_q\), since
\(q^{-1}=J_2(h)p^{-1}J_1(h)^{-1}\).

For a connected complex Lie group \(L\), the character
\(\det\Ad_L:L\to\C^\times\) has differential
\(x\mapsto\tr(\ad_{\operatorname{Lie}(L)}x)\).
If this differential is zero, the character is locally constant,
hence identically one.  This applies to \(G\) by unimodularity,
and to \(N\) because a semisimple Lie algebra is perfect and the trace
of a commutator is zero.  Taking determinants in
\eqref{eq:covariance} yields
\[
 f(A_h(p))=f(p).
\]
Consequently \(f=f(e)\ne0\) on the open identity orbit.
The identity theorem on the connected complex manifold \(N\) gives
\(f\equiv f(e)\) on all of \(N\).
Every \(B_p\) is therefore invertible, so every orbit is open.
Connectedness of \(N\) implies that there is only one orbit:
the action \eqref{eq:action} is transitive.

Let \(H=\{h\in G:A_h(e)=e\}\).
It is a closed subgroup with Lie algebra
\(\ker B_e=0\), hence is discrete.
The orbit map identifies \(G/H\) with \(N\) as a manifold
\cite[Proposition~4.12 and Corollary~4.16]{Etingof}.
For completeness, \(G\to G/H\) is a covering: choose an identity
neighborhood \(U\) with \(U^{-1}U\cap H=\{e\}\); the sets \(Uh\),
\(h\in H\), are disjoint sheets over the image of \(U\), and left
translations give such neighborhoods everywhere.
Because \(G\) is connected and \(N\) is simply connected, this
covering has one sheet \cite[Proposition~1.32]{Hatcher}.
Thus \(H=\{e\}\), and
\begin{equation}\label{eq:diffeo}
 G\ \text{is diffeomorphic to}\ N .
\end{equation}
The orbit map need not be a group homomorphism; no group or
Lie algebra isomorphism is inferred from \eqref{eq:diffeo}.

Take a Levi decomposition \(\g=\rad\rtimes\s\), where
\(\rad\) is the solvable radical and \(\s\) is semisimple.
Let \(R,S\) be their simply connected integrations.
The Levi action integrates on \(R\); the simply connected group
\(R\rtimes S\) has Lie algebra \(\g\).
By uniqueness of integration, its manifold \(R\times S\) is the
manifold of \(G\)
\cite[Theorem~48.9 and Corollary~9.14]{Etingof}.
This construction imposes no closed-image assumption on a
previously chosen Levi subgroup.  Lemma~\ref{lem:solvable}
makes the \(R\) factor contractible.

Polar decomposition retracts a simply connected complex semisimple
group onto a compact real form
\cite[Corollary~43.5, Theorem~43.7, Corollary~43.8 and the
following covering discussion]{Etingof}.
Let these compact groups be \(K_N\) for \(N\) and \(K_S\) for \(S\).
Their real dimensions are respectively \(d\) and
\(s=\dim_{\C}\s\), rather than \(2d\) and \(2s\).
The compact connected Lie group \(K_N\) is orientable: left
translation of a basis at its identity gives a global frame.
Its fundamental class therefore gives, in ordinary singular homology
\cite[Theorem~3.26(a)]{Hatcher},
\[
 H_d(N;\Z)\cong H_d(K_N;\Z)\cong\Z.
\]
If \(\rad\ne0\), then \(s<d\).
The product description of \(G\), contractibility of \(R\), and polar
decomposition for \(S\) give
\[
 H_d(G;\Z)\cong H_d(K_S;\Z)=0,
\]
since \(K_S\) is a compact \(s\)-manifold
\cite[Theorem~3.26(c)]{Hatcher}.
This contradicts \eqref{eq:diffeo}.
Hence \(\rad=0\), proving that \(\g\) is semisimple.
\end{proof}

\section{The post-Lie consequence}

Let \(\g=(V,[\, ,\,])\) and \(\n=(V,\{\, ,\,\})\).
Our convention for a post-Lie product \(x\cdot y\) is
\begin{align}
 x\cdot y-y\cdot x&=[x,y]-\{x,y\},\label{eq:PL1}\\
 [x,y]\cdot z&=x\cdot(y\cdot z)-y\cdot(x\cdot z),\label{eq:PL2}\\
 x\cdot\{y,z\}&=\{x\cdot y,z\}+\{y,x\cdot z\}.\label{eq:PL3}
\end{align}
Thus \(L(x)y=x\cdot y\) defines a homomorphism
\(L:\g\to\Der(\n)\).

\begin{corollary}\label{cor:postlie}
If a finite-dimensional complex pair \((\g,\n)\) admits a post-Lie
product, \(\g\) is unimodular and \(\n\) is semisimple, then
\(\g\) is semisimple.  In particular, a perfect nonsemisimple complex
source and a semisimple target admit no post-Lie product.
\end{corollary}

\begin{proof}
For semisimple \(\n\), every derivation is inner and its center is
zero.  To recall the innerness argument, let \(\kappa\) be its
nondegenerate Killing form and \(D\in\Der(\n)\).
Choose \(a\in\n\) such that
\(\kappa(a,u)=\tr(D\ad_{\n}u)\).
Using \([D,\ad_{\n}u]=\ad_{\n}(Du)\), cyclicity of trace and
invariance of \(\kappa\), we obtain
\[
 \kappa(Du,v)
 =\tr\bigl(D[\ad_{\n}u,\ad_{\n}v]\bigr)
 =\kappa(a,\{u,v\})=\kappa(\{a,u\},v).
\]
Hence \(D=\ad_{\n}a\).
The map \(\ad_{\n}:\n\to\Der(\n)\) is therefore a linear
isomorphism.  In particular there is a unique linear map
\[
 R=\ad_{\n}^{-1}L:V\longrightarrow V,\qquad
 x\cdot y=\{Rx,y\}.
\]
Equations \eqref{eq:PL2}--\eqref{eq:PL3} and injectivity of
\(\ad_{\n}\) give
\[
 R[x,y]=\{Rx,Ry\}.
\]
Equation \eqref{eq:PL1} gives the signs explicitly:
\[
 [x,y]=\{Rx,y\}+\{x,Ry\}+\{x,y\}.
\]
Expanding \(\{(R+I)x,(R+I)y\}\) now shows that
\[
 j_1=R+I,\qquad j_2=R
\]
are homomorphisms \(\g\to\n\), with difference \(I\).
Theorem~\ref{thm:pair} applies.
Finally the functional \(x\mapsto\tr(\ad_{\g}x)\) vanishes on
\([\g,\g]\), since
\(\ad_{\g}[x,y]=[\ad_{\g}x,\ad_{\g}y]\).
A perfect source is consequently unimodular.
\end{proof}

\begin{remark}
The zero-dimensional algebra is semisimple by convention, so it gives
no exception to the nonsemisimple-source assertion.
Unimodularity is essential to Theorem~\ref{thm:pair}.
For example, let \(\g\) have basis \(H,E,Z\), with
\([H,E]=2E\) and \(Z\) central.
In \(\n=\mathfrak{sl}_2(\C)\), use the standard basis \(H,E,F\).
The maps
\[
 j_1(H)=H,\quad j_1(E)=E,\quad j_1(Z)=0,\qquad
 j_2(H)=j_2(E)=0,\quad j_2(Z)=F
\]
are homomorphisms with invertible difference, but
\(\tr(\ad_{\g}H)=2\).
The proof makes no claim for a general real form, positive
characteristic, infinite dimension, or a merely reductive target.
It assumes neither algebraicity nor completeness of a locally
defined action, and no bijectivity of a solvable exponential map.
\end{remark}

\section*{Verification and authorship disclosure}

The accompanying archival checker reports 2,896 exact finite
diagnostic conditions in each run.  Those calculations are auxiliary;
the proof above does not infer integration, transitivity, covering
theory or homology from finite testing.
AI systems were used extensively in solving, checking and writing
this work.  Separate AI audits examined the algebraic, global-action
and topological arguments.  This is an unrefereed preprint: no
completed human peer review or formal machine verification is claimed.

\begin{thebibliography}{9}

\bibitem{Baues}
O.~Baues,
\emph{Prehomogeneous affine representations and flat
pseudo-Riemannian manifolds},
\href{https://arxiv.org/abs/0809.0824v1}{arXiv:0809.0824v1} (2008).

\bibitem{BD}
D.~Burde and K.~Dekimpe,
\emph{Post-Lie algebra structures and generalized derivations of
semisimple Lie algebras},
Mosc. Math. J. \textbf{13} (2013), no.~1, 1--18.
\href{https://homepage.univie.ac.at/dietrich.burde/papers/burde_44_postlie_der.pdf}
{Author manuscript dated 10 September 2012}.

\bibitem{BDV}
D.~Burde, K.~Dekimpe and K.~Vercammen,
\emph{Affine actions on Lie groups and post-Lie algebra structures},
Linear Algebra Appl. \textbf{437} (2012), 1250--1263;
\href{https://arxiv.org/abs/1109.0251}{arXiv:1109.0251}.

\bibitem{BDM}
D.~Burde, K.~Dekimpe and M.~Monadjem,
\emph{Post-Lie algebra structures for perfect Lie algebras},
Comm. Algebra \textbf{52} (2024), no.~10, 4255--4267,
\href{https://doi.org/10.1080/00927872.2024.2344638}
{doi:10.1080/00927872.2024.2344638}.
Theorem numbering follows the
\href{https://homepage.univie.ac.at/dietrich.burde/papers/burde_78_perfect.pdf}
{author manuscript dated 8 March 2024};
preprint \href{https://arxiv.org/abs/2311.08985}{arXiv:2311.08985} (2023).

\bibitem{OWR}
D.~Burde (joint work with K.~Dekimpe),
\emph{Pre-Lie and post-Lie algebra structures},
in \emph{Mini-Workshop: Poisson and Poisson-type algebras},
Oberwolfach Report 46/2023, 2688--2690,
\href{https://doi.org/10.4171/OWR/2023/46}
{doi:10.4171/OWR/2023/46}; conjecture at p.~2689.

\bibitem{Etingof}
P.~Etingof,
\emph{Lie groups and Lie algebras},
\href{https://ocw.mit.edu/courses/18-755-lie-groups-and-lie-algebras-ii-spring-2024/mit18_755_s24_lec_full.pdf}
{MIT 18.755 lecture notes, Spring 2024}.

\bibitem{Hatcher}
A.~Hatcher,
\emph{Algebraic Topology},
Cambridge University Press, 2002;
\href{https://pi.math.cornell.edu/~hatcher/AT/AT.pdf}{author's text}.

\end{thebibliography}
\end{document}
